Bei jedem bestimmten Zinsfuß kann eine Kapitalzunahme von 5000 Thl. die Pensionserhöhung um 10 Thl. nur für eine bestimmte Zahl von Pensionierten decken; bei einem Zinsfuß von 3 Prozent\footnote{Als der Vorschlag zu der fraglichen Regelung im Juli 1794 eingebracht wurde, rechnete man in unzweideutigen Worten nur mit dieser Verzinsung und nicht mit 4 Prozent. Danach ist also die folgende Stelle im Aufsatz des Herrn Universitätsrats O. vom 15. Dezember 1844, S. 5:
\begin{quote}\raggedright
‘Bei der im Jahr 1794 getroffenen Bestimmung .\,.\,.\,.\,. ist wahrscheinlich eine Berechnung zugrunde gelegt worden,\par
‘nach der die Zinsen des auf 5000 Thl. erhöhten Fonds bei 4 Prozent 200 Thl. betra-\par
‘gen, davon \(\tfrac14\) wieder zum Kapital zu schlagen, \(\tfrac34\) aber unter die Witwen zu verteilen\par
‘seien, was dann bei 15 Witwen für jede 10 Thl. betragen würde’
\end{quote}
zu berichtigen. Von einer solchen Berechnung, von 4 Prozent Zinsen und von der Zurücklegung eines Viertels davon findet sich in den bald näher zu betrachtenden Verhandlungen von 1794 \textit{gar nichts}.} für höchstens
\SourceLeaf{136}
15, bei einem Zinsfuß von \(3\tfrac12\) Prozent für höchstens 17 und bei einem Zinsfuß von 4 Prozent für höchstens 20 Pensionierte; der Zinsertrag jener Kapitalvermehrung wird, wenn diese Zahlen erreicht sind, dadurch vollständig aufgezehrt. Steigt also die Zahl der zu Pensionierenden über die jeweiligen Grenzen, so ist zur Bestreitung aller übrigen Pensionen nichts vorhanden außer 1) dem Überschuss, den die jährlichen Einnahmen der Kasse vor den Erhöhungen des Kapitals und der Pensionssätze erbracht hatten, oder vielmehr erbracht hätten, wenn die jeweilige festgelegte Pensionszahl (15; 17; 20) damals bestanden hätte\footnote{In Wirklichkeit hatten die Einnahmen damals (1794) schon für 18 Pensionierungen nicht ganz ausgereicht, obwohl zu jener Zeit noch ein jährlicher Zuschuss von 150 Thl. aus der Kirchenkasse gezahlt wurde, der später aufgehört hat. Die obige Überschlagsrechnung erfährt jedoch eine Veränderung, weil die Progressionsregel vor allem infolge von Unregelmäßigkeiten während der westphälischen Regierung nicht genau befolgt wurde. Hätte man sich ganz streng daran gehalten, so hätten die Pensionssätze seit 1815 stets schon bei einem geringeren Kapitalstand erhöht werden müssen, als es tatsächlich geschah, — und das jetzige Kapitalvermögen wäre um vielleicht 5000 Thl. ärmer.}. 2) Der etwaigen seit jener Zeit bewirkten Steigerung der Apothekenpacht. 3) Dem vergrößerten Ertrag der Mitgliederbeiträge infolge ihrer gewachsenen Zahl. Diese unsicheren und schwachen Hilfsquellen wären jedoch bei weiterem Überschreiten jener festgelegten Zahlen bald erschöpft, und zwar umso schneller, je höher der Pensionssatz selbst bereits angewachsen ist. Diese letzte Bemerkung ist insofern von großer Bedeutung, als daraus mit größter Klarheit hervorgeht, dass die Gefahr aus einer konsequenten Befolgung der zweiten Interpretation umso größer wird, je mehr Pensionserhöhungen bis zum Überschreiten der festgelegten Zahl bereits stattgefunden haben.

Wenn oben bemerkt wurde, dass im gedruckten Regulativ keinerlei Vorbehalt zu finden ist, mit dem ein früher oder später aus dieser Ursache folgendes Ausbluten der Kasse aufgehalten werden könnte, so darf ich nicht verschweigen, dass in die Quittungsformulare, mit denen sich die Witwen ihre Pensionen auszahlen lassen, der Vorbehalt aufgenommen wurde, die Pension sei auf .\,.\,.\,.\,. Rthl. für jetzt und \textit{solange die Umstände der Kasse es gestatten werden} festgesetzt. Wahrscheinlich kennen nur wenige Mitglieder der Witwenkasse diese Klausel; mir selbst ist sie jedenfalls erst ganz vor Kurzem bekannt geworden, obwohl ich 38 Jahre Teilnehmer gewesen bin. Sie ist jedoch nicht dazu bestimmt, wie ich anfangs vermutete, möglichen Gefahren aus der Progressionsregelung vorzubeugen; denn sie steht mit genau denselben Worten auch schon in den gedruckten Quittungsformularen vor 1794. Solche ganz allgemein gefassten Vorbehalte mögen in mancher Hinsicht ihr Gutes haben. Im gewöhnlichen Lauf der Dinge bleiben sie aber, wenn keine außerordentlichen Anlässe eintreten, so lange ohne Anwendung, bis die höchste Not dazu zwingt. Es ist kein Entscheidungskriterium, kein Maßstab angegeben, an dem man erkennen kann, ob die Umstände der Kasse die Zahlung gestatten oder nicht gestatten. Wartet man so lange, bis man schon wiederholt gezwungen ist, zur Zahlung der Pensionen auch das Kapital anzugreifen, so hat man sehr wahrscheinlich bereits viel zu lange gewartet; dann wird sich bewähren, was [S. 128, Z. 12] bemerkt wurde. Jedenfalls besteht für die Bewahrung des Rufs der Kasse als solide Einrichtung ein großer Unterschied zwischen einem Rückgang des Pensionssatzes infolge einer ganz bestimmten, feststehenden Regel (wie bei der ersten Interpretation) und einer Kürzung der Pensionen, zu der sich die Kasse schließlich gezwungen sieht, weil sie die Verpflichtungen nicht mehr erfüllen kann, die sie nach dem
\SourceLeaf{137}
Statut (in der zweiten Interpretation) eigentlich \textit{unbedingt} übernommen hat.

